% Generated draft: requires mathematical and source review.
\documentclass[11pt]{article}
\usepackage[margin=25mm]{geometry}
\usepackage{fontspec}
\setmainfont{Latin Modern Roman}
\usepackage{amsmath,amssymb,mathtools,unicode-math}
\setmathfont{Latin Modern Math}
\usepackage{enumitem,needspace,xurl,hyperref}
\setlength{\parindent}{0pt}
\setlength{\parskip}{5pt}
\setlength{\emergencystretch}{3em}
\newcommand{\N}{\mathbb{N}}
\newcommand{\Z}{\mathbb{Z}}
\newcommand{\Q}{\mathbb{Q}}
\newcommand{\R}{\mathbb{R}}
\newcommand{\C}{\mathbb{C}}
\newcommand{\F}{\mathbb{F}}
\newcommand{\A}{\mathbb{A}}
\newcommand{\PP}{\mathbb{P}}
\newcommand{\E}{\mathbb{E}}
\newcommand{\eps}{\varepsilon}
\newcommand{\dd}{\,\mathrm{d}}
\begin{document}
\section*{UnsolvedMath ID: 20002720}
\textbf{Source problem code:} catalogue entry
\subsection{Definitions and mathematical statement}
\textbf{Problem 5.2 (A. Nica).} Describe the unit group for multivariable boxed convolution, and find a multivariable analogue of the \(S\)-transform/free Fourier transform. Let \(\mathcal S_s=\C_0\langle\!\langle z_1,\ldots,z_s\rangle\!\rangle\), \(s\geq1\), be formal noncommutative series with zero constant coefficient. For a word \(w=(i_1,\ldots,i_n)\), write \(f_w\) for its coefficient. If \(\pi\in\operatorname{NC}(n)\), set \(f_{w|\pi}=\prod_{B\in\pi}f_{w|B}\), with each subword read in increasing position order. Using the right Kreweras complement \(K(\pi)\), define
\[
 (f\boxtimes g)_w=\sum_{\pi\in\operatorname{NC}(n)}
 f_{w|\pi}g_{w|K(\pi)}.
\]
The complement convention is fixed by the usual maximal noncrossing complement on interleaved unbarred/barred positions. The source's corrupted ``?s'' denotes this operation. It is associative with unit \(\Delta_s=z_1+\cdots+z_s\). For tuples in a unital probability algebra, freeness is between the algebras generated by the two whole tuples; then the source relation is \(R_a\boxtimes R_b=R_{(a_1b_1,\ldots,a_sb_s)}\). No internal independence within each tuple is assumed.

\textbf{Reported units and central splitting.} A formal series is invertible exactly when
\[
 f_{(i)}\ne0\qquad(1\leq i\leq s).
\]
Its inverse is determined recursively by word length. Explicitly, if \(a_i=f_{(i)}\), \(a_w=\prod_ja_{i_j}\), the right-inverse coefficients start at \(g_{(i)}=a_i^{-1}\), and for \(|w|\geq2\) satisfy
\[
 g_w=-a_w^{-1}\left[f_wa_w^{-1}
 +\sum_{\pi\ne0_n,1_n}f_{w|\pi}g_{w|K(\pi)}\right].
\]
All coefficients in the sum have smaller word length; a left recursion gives the same group inverse.

For \(a\in(\C^*)^s\), let \(T_a=\sum_i a_iz_i\). The aid reports the coefficientwise central action
\[
 (T_a\boxtimes f)_w=(f\boxtimes T_a)_w=a_wf_w.
\]
If \(U_s\) is the subgroup with all linear coefficients one, every unit has the unique factorization \(f=T_a\boxtimes u\), \(u\in U_s\), hence the canonical direct product
\[
 \mathcal S_s^\times\cong(\C^*)^s\times U_s.
\]
This is a formal scalar-coefficient statement, not an analytic or operator-valued transform theorem.

For \(p\geq2\), define \(F^pU_s\) by vanishing of all nonlinear coefficients of degrees less than \(p\). The reported commutator estimate is
\[
 [F^pU_s,F^qU_s]\subseteq F^{p+q-1}U_s\qquad(p,q\geq2).
\]
Finite degree truncations are nilpotent unipotent algebraic groups, and \(U_s\) is their degree-complete inverse limit. The aid reports a first nonzero graded commutator in degree three, proving noncommutativity for \(s\geq2\); its complete bracket formula is not included in the supplied summary. Consequently no injective group homomorphism of this full unit group into an abelian multiplicative target can provide a faithful multivariable transform. Noncommutative target constructions are not excluded or completed by this obstruction.

\textbf{One-variable source remarks and target correction.} For \(s=1\), boxed convolution is commutative, and the source cites Nica--Speicher's free Fourier isomorphism \(F\), satisfying \(F(R_a)=S_a\) when \(\varphi(a)\ne0\). With the above cumulant-series convention, the usual expression is \(F(f)(z)=f^{\langle-1\rangle}(z)/z\), using compositional inverse. Its multiplicative target must be \(\C[[z]]^\times\), the series with nonzero constant coefficient. The literal source target \(\C_0[[z]]\) under ordinary multiplication has no unit and cannot have the claimed unit group. This is an explicit notation defect, not a silent redefinition. The requested multivariable \(S\)-transform remains broader than the supplied group-structure result.
\subsection{Short English statement}
The aid describes boxed-convolution units by nonzero linear coefficients and splits their group into a central torus and a normalized inverse limit of nilpotent groups. Multivariable noncommutativity blocks faithful transforms into abelian multiplicative targets. The one-variable remark also needs a target correction: ordinary multiplicative units have nonzero constant term.
\subsection{Sources}
\begin{itemize}
\item Local UnsolvedMath registry, permanent ID \(20002720\), catalogue entry, \texttt{aim-probability-notes.json}, \texttt{notes[161]}, workshop \emph{Braid groups, clusters and free probability}, Problem 5.2, A. Nica: \url{https://aimath.org/WWN/braidgroups/braidgroups.pdf}.
\item Original references [39], Alexandru Nica and Roland Speicher, multivariable boxed convolution/noncrossing partitions; [38], their free Fourier transform; the remarks also refer to Voiculescu's \(S\)-transform. Full bibliography entries are not supplied.
\item Supplied research summary on inverse recursions, central torus, degree filtration and first graded bracket. Attempt 1: partial\_result, partially\_solved, rigor ``proved'', low novelty confidence. The machine-generated aid is not peer reviewed and requires independent verification before citation. The supplied verification identifies Problem 5.2; this draft has not visited the PDF. The extracted 11 is a page marker.
\end{itemize}
\end{document}
